% Part: proof-theory
% Chapter: sequent-calculus
% Section: rules-proofs

\documentclass[../../../include/open-logic-section]{subfiles}

\begin{document}

\olfileid{pt}{seq}{rp}

\olsection{القواعد و\usetoken{P}{proof}}

نبدأ بوضع بعض التعريفات وتثبيت بعض المصطلحات.

لننظر، مثلًا، في قاعدتي~\Log{G3c} اللتين تتيحان~!!{proving} متتاليات
تحتوي~$!A\land !B$ انطلاقًا من متتاليات تحتوي~$!A$ و$!B$:
\[
\Axiom$ !A, !B, \Gamma \fCenter \Delta$
\RightLabel{\LeftR{\land}}
\UnaryInf$ !A \land !B, \Gamma \fCenter \Delta$
\DisplayProof
\qquad
\Axiom$\Gamma \fCenter \Delta, !A$
\Axiom$ \Gamma \fCenter \Delta, !B$
\RightLabel{\RightR{\land}}
\BinaryInf$ \Gamma \fCenter \Delta, !A \land !B $
\DisplayProof
\]
تسمى المتتاليات الواقعة فوق الخط \emph{مقدمات}، وتسمى المتتالية الواقعة
تحته \emph{نتيجة}. وتسمى~!!{formula}s الموجودة في~$\Gamma$ و$\Delta$
في كل قاعدة \emph{السياق} أو~!!{formula}s \emph{الجانبية}. أما
!!{formula} في النتيجة التي ليست جزءًا من السياق فتسمى~!!{formula}
\emph{الرئيسة}، وتسمى~!!{formula}s الموجودة في المقدمات لا في السياق
!!{formula}s \emph{الفعالة} (أو \emph{المساعدة}).

قواعد الكميات أعقد قليلًا. فقواعد~$\lforall$ في~\Log{G1c} مثلًا هي:
\[\Axiom$ !A(t), \Gamma \fCenter \Delta$
\RightLabel{\LeftR{\lforall}}
\UnaryInf$ \lforall[x][!A(x)],\Gamma \fCenter \Delta$
\DisplayProof
\qquad
\Axiom$ \Gamma \fCenter \Delta, !A(a) $
\RightLabel{\RightR{\lforall}}
\UnaryInf$ \Gamma \fCenter \Delta, \lforall[x][!A(x)]$
\DisplayProof
\]
في قاعدة~\LeftR{\lforall} تكون~!!{formula} الفعالة هي~$!A(t)$، حيث
$t$ حد مغلق، أي حد بلا متغيرات. ويكون مثل هذا الاستدلال صحيحًا---أي
مطابقًا لمخطط قاعدة~\LeftR{\lforall}---إذا وجدت~!!{formula}~$!A$
ومتغير~$x$ وحد مغلق~$t$ بحيث تكون~!!{formula} الرئيسة في النتيجة
$\lforall[x][!A]$ وتكون~!!{formula} الفعالة في المقدمة
$\Subst{!A}{t}{x}$. وتعمل قاعدة~\RightR{\lforall} على النحو نفسه، إلا
أن~!!{formula} الفعالة في المقدمة يجب أن تكون~$\Subst{!A}{a}{x}$، حيث~$a$
!!a{constant}. ويسمى هذا الثابت \emph{المتغير المميَّز} لاستدلال
\RightR{\lforall}. ولا يصح الاستدلال إلا إذا لم تحتوي المتتالية السفلى
$a$؛ أي إذا لم يقع~$a$ في~$\Gamma$ أو~$\Delta$ أو~$!A$. وهذا هو
\emph{شرط المتغير المميَّز}.

تتكون أنظمة المتتاليات مثل~\Log{G1c} و\Log{G3c} من مجموعة قواعد محددة
مخططيًا كهذه، ومعها بعض المتتاليات المحددة مخططيًا أيضًا التي تعد
بديهيات. أما~!!^{proof}s في مثل هذا النظام فهي أشجار منتهية من
المتتاليات، تُعرَّف استقرائيًا من البديهيات باستعمال قواعد الاستدلال. فكل
ورقة بديهية، وكل متتالية أخرى تتبع من المتتاليات الواقعة فوقها مباشرة
بإحدى قواعد الاستدلال في النظام.

\begin{defn}\ollabel{defn:proof}
تُعرَّف~\emph{!!^{proof}s} في نظام متتاليات استقرائيًا بوصفها أشجارًا
منتهية من المتتاليات كما يأتي:
\begin{enumerate}
  \item\ollabel{defn:prf:axiom} كل بديهية~!!a{proof}.
  \item\ollabel{defn:prf:one-prem} إذا كان~$\pi_1$~!!a{proof} ينتهي
  بمتتالية~$S_1$، وكانت~$S$ متتالية تتبع من~$S_1$ بقاعدة أحادية
  المقدمة~$R$، فإن
  \[
  \AxiomC{}
  \RightLabel{$\pi_1$}
  \DeduceC{$S_1$}
  \RightLabel{$R$}
  \UnaryInfC{$S$}
  \DisplayProof
  \]
  هو~!!a{proof}.
  \item\ollabel{defn:prf:two-prem} إذا كان~$\pi_1$ و$\pi_2$~!!{proof}s
  ينتهيان بالمتتاليتين~$S_1$ و$S_2$، وكانت~$S$ متتالية تتبع من~$S_1$
  و~$S_2$ بقاعدة ثنائية المقدمة~$R$، فإن
  \[
  \AxiomC{}
  \RightLabel{$\pi_1$}
  \DeduceC{$S_1$}
  \AxiomC{}
  \RightLabel{$\pi_2$}
  \DeduceC{$S_2$}
  \RightLabel{$R$}
  \BinaryInfC{$S$}
  \DisplayProof
  \]
  هو~!!a{proof}.
\end{enumerate}
\end{defn}

يعرّف~\Olref{defn:proof}~!!{proof}s بوصفها أشجارًا من المتتاليات. ونتصور
هذه الأشجار نامية «إلى أعلى». والعقد العليا (الأوراق) بديهيات وتسمى
\emph{متتاليات ابتدائية}. أما المتتالية السفلى فتسمى \emph{المتتالية
النهائية}. وإذا كان~$\pi$~!!a{proof} متتاليته النهائية~$S$، نكتب أيضًا
$\pi\Proves S$. وإذا وجد~!!a{proof} لمتتالية~$S$، كتبنا~$\Proves S$،
وقد نشير إلى النظام الذي يوجد فيه البرهان إلى يسار رمز~$\Proves$؛ مثلًا
$\Log{G1c}\Proves !A,!B\Sequent !A\land !B$. وكل متتالية في
!!a{proof}، عدا المتتاليات الابتدائية، هي \emph{نتيجة} استدلال بقاعدة
ما~$R$. وتسمى المتتالية أو المتتاليتان الواقعتان مباشرة فوق متتالية في
!!a{proof} (أي أبوا العقدة) \emph{مقدمات} الاستدلال.

تسمى~!!{proof}s المستعملة في البناء الاستقرائي لـ!!a{proof}
\emph{!!{proof}s جزئية} له. وتسمى~!!{proof}s الجزئية المنتهية بمقدمات
استدلال \emph{!!{proof}s جزئية مباشرة} لذلك الاستدلال.

كثيرًا ما يلزم قياس تعقيد~!!{proof}s بعدد طبيعي. ويمكن ببساطة عدّ
المتتاليات في~!!a{proof}، لكن المقياس الأنفع غالبًا هو~!!{height}
!!a{proof}: وهو أكبر عدد من خطوات الاستدلال، أو الحواف، في فرع من
المتتالية النهائية إلى متتالية ابتدائية. ويمكن أيضًا إعطاؤه تعريفًا
استقرائيًا.

\begin{defn}
يُعرَّف~\emph{!!{height}}~$\pheight{\pi}$ لـ!!a{proof}~$\pi$ استقرائيًا
كما يأتي:
\begin{enumerate}
  \item إذا اقتصر~$\pi$ على بديهية واحدة، كان~$\pheight{\pi}=0$.
  \item إذا انتهى~$\pi$ باستدلال ذي مقدمة واحدة وكان~$\pi_1$ هو
  الـ!!{proof} الجزئي المباشر، كان~$\pheight{\pi}=\pheight{\pi_1}+1$.
  \item إذا انتهى~$\pi$ باستدلال ذي مقدمتين وكان~$\pi_1$ و$\pi_2$
  الـ!!{proof}s الجزئية المباشرة، كان
  $\pheight{\pi}=\max(\pheight{\pi_1},\pheight{\pi_2})+1$.
\end{enumerate}
إذا كان للمتتالية~$S$~!!a{proof} ارتفاعه~$\le n$، كتبنا~$\Proves[n]S$.
\end{defn}

\begin{ex}
في~\Log{G3c} تكون كل متتالية من الصورة~$!C,\Gamma\Sequent\Delta,!C$
بديهية، على أن تكون~$!C$ ذرية. ولتكن~$!D$ و$!E$ جملتين ذريتين. عندئذ
تكون~$!D,!E\Sequent !D$ و$!D,!E\Sequent !E$ حالتين من البديهية
$!C,\Gamma\Sequent\Delta,!C$. فمثلًا، إذا أخذنا~$!C\ident !E$ و
$\Gamma=\{!D\}$ و$\Delta=\emptyset$، حصلنا بالضبط على
$!D,!E\Sequent !E$. (تذكر أننا نتعامل مع متعددات مجموعات؛ ولذلك فإن
$!D,!E\Sequent !E$ و$!E,!D\Sequent !E$ هما المتتالية نفسها.)

ولأن~$!D,!E\Sequent !D$ حالة من بديهيات~\Log{G3c}، فإنه يعد وحده
!!a{proof} في~\Log{G3c}، وكذلك~$!D,!E\Sequent !E$ بحسب
\olref{defn:proof}\olref{defn:prf:axiom}. ولنسميهما~$\pi_1$ و$\pi_2$.
وبحسب~\olref{defn:prf:two-prem}، يمكننا أخذ هذين~!!{proof}s وتوليد واحد
جديد~($\pi_3$) بقاعدة~\RightR{\land} ثنائية المقدمة:
\[
\Axiom$!D, !E \fCenter !D$
\Axiom$!D, !E \fCenter !E$
\RightLabel{\RightR{\land}}
\BinaryInf$!D, !E  \fCenter !D \land !E $
\DisplayProof
\]
ومن~$\pi_3$ نحصل بدوره على~!!{proof}~$\pi_4$:
\[
\Axiom$!D, !E \fCenter !D$
\Axiom$!D, !E \fCenter !E$
\RightLabel{\RightR{\land}}
\insertBetweenHyps{\hspace{5ex}}
\BinaryInf$!D, !E  \fCenter !D \land !E $
\LeftSubproofLabel{$\pi_3$}
\RightLabel{\LeftR{\land}}
\UnaryInf$!D \land !E  \fCenter !D \land !E$
\RightSubproofLabel{$\pi_4$}
\DisplayProof
\]
با استعمال قاعدة~\LeftR{\land} أحادية المقدمة (وبحسب
\olref{defn:proof}\olref{defn:prf:one-prem}). والـ!!{proof} الجزئي
المباشر للاستدلال الأخير في~$\pi_4$ هو~$\pi_3$. والـ!!{proof}s الجزئية
المباشرة لاستدلال~\RightR{\land} هي~$\pi_1$ و$\pi_2$. وتشمل
الـ!!{proof}s الجزئية لـ$\pi_4$ كلًا من~$\pi_1$ و$\pi_2$ و$\pi_3$
و$\pi_4$ نفسه. ولدينا~$\pheight{\pi_1}=\pheight{\pi_2}=0$ و
$\pheight{\pi_3}=1$ و$\pheight{\pi_4}=2$. وقد أثبتنا أن
$\Log{G3c}\Proves[2] !D\land !E\Sequent !D\land !E$ لأي~$!D$ و$!E$
ذريتين.
\end{ex}

\end{document}
